\par\SourceLeaf{424}%
\noindent\hspace*{1.2em}The proof is easily given as follows. Let
\[
AL=t,\quad AL'=t',\quad AL''=t'',\quad AL'''=t''';
\]
Then
\[
\begin{aligned}
\cos LL'&=\cos t\cos t'+\sin t\sin t'\cos A\\
\cos L''L'''&=\cos t''\cos t'''+\sin t''\sin t'''\cos A\\
\cos LL''&=\cos t\cos t''+\sin t\sin t''\cos A\\
\cos L'L'''&=\cos t'\cos t'''+\sin t'\sin t'''\cos A.
\end{aligned}
\]
Hence
\[
\begin{aligned}
&\cos LL'\cos L''L'''-\cos LL''\cos L'L'''\\
&=\cos A\{\cos t\cos t'\sin t''\sin t'''+\cos t''\cos t'''\sin t\sin t'\\
&\qquad\qquad-\cos t\cos t''\sin t'\sin t'''-\cos t'\cos t'''\sin t\sin t''\}\\
&=\cos A(\cos t\sin t'''-\cos t'''\sin t)(\cos t'\sin t''-\cos t''\sin t')\\
&=\cos A\sin(t'''-t)\sin(t''-t')\\
&=\cos A\sin LL'''.\sin L'L''.
\end{aligned}
\]

Since each of the two great circles extends from $A$ in two opposite directions, two angles that add up to $180^{0}$ are formed there. Our analysis shows, however, that one must choose the rays that run in the sense from $L$ to $L'''$ and from $L'$ to $L''$.

In place of the angle $A$ we may also take the distance from the pole of the great circle $LL'''$ to the pole of the great circle $L'L''$. Since each great circle has two poles, we see that we must join the poles about which the respective great circles run from $L$ to $L'''$ and from $L'$ to $L''$ in the same sense.

We leave to the reader the treatment of the special case in which one or both of the arcs $LL'''$ and $L'L''$ measure $90^{0}$.

6) Another useful proposition follows from the following analysis. Let $L,L',L''$ be three points on the sphere and let
\[
\begin{aligned}
\cos L(1)&=x,&\quad\cos L(2)&=y,&\quad\cos L(3)&=z,\\
\cos L'(1)&=x',&\cos L'(2)&=y',&\cos L'(3)&=z',\\
\cos L''(1)&=x'',&\cos L''(2)&=y'',&\cos L''(3)&=z''.
\end{aligned}
\]
